% ============================================================================
\chapter{Configuration}
\label{ch:configuration}
% ============================================================================

All settings live in a single file, \code{config.toml}, read once at start-up from
the application's working directory. \textbf{Changes take effect after restarting the
application} (\code{sudo systemctl restart deltax2} on a kiosk installation).

If the file is missing or contains a syntax error, the application prints an error
and exits; on a kiosk installation check \code{journalctl -u deltax2} for the reason.

\section{Serial connection --- \code{[serial]}}

\begin{lstlisting}[language=TOML]
[serial]
port = "/dev/ttyUSB0"   # robot serial device
baud_rate = 115200      # must match the robot firmware
\end{lstlisting}

\begin{center}
\begin{tabularx}{\textwidth}{@{}llX@{}}
  \toprule
  \textbf{Key} & \textbf{Type} & \textbf{Meaning} \\
  \midrule
  \code{port} & string & Serial device path. Typically \code{/dev/ttyUSB0} or
    \code{/dev/ttyACM0} on Linux. Use the \textbf{List com} button
    (Section~\ref{sec:systemscreen}) or \code{ls /dev/tty*} to find it. \\
  \code{baud\_rate} & integer & Communication speed; the Delta~X~2 firmware uses
    115\,200 (Appendix~\ref{app:gcode}). \\
  \bottomrule
\end{tabularx}
\end{center}

\begin{tip}
USB device names can change between reboots when several adapters are connected.
For a robust kiosk setup, use a stable symlink from
\code{/dev/serial/by-id/\dots} as the \code{port} value.
\end{tip}

\section{User interface --- \code{[ui]}}

\begin{lstlisting}[language=TOML]
[ui]
kiosk_mode = false   # true: borderless full-screen for the Pi touch screen
\end{lstlisting}

\section{Safety limits --- \code{[robot]}}
\label{sec:limits}

\begin{lstlisting}[language=TOML]
[robot]
limit_min = { x = -160.0, y = -160.0, z = -200.0 }
limit_max = { x =  160.0, y =  160.0, z =    0.0 }
\end{lstlisting}

Every jog and seeding movement is checked against these boundaries \emph{before} any
command is sent to the hardware; a move that would leave the box is refused with the
status message \emph{Movement out of safety limits}. The values must lie within the
physical workspace of the robot, described in Section~\ref{sec:workspace}.

\begin{warning}
The coordinate origin is established by homing (\code{G28}): after homing, the head
is at the \emph{top centre} of the working volume, at $X{=}0$, $Y{=}0$, $Z{=}0$.
The Z axis points \emph{upwards}: all working positions below the home position have
\textbf{negative Z}. Consequently \code{limit\_max.z} should normally stay at
\code{0.0} and \code{limit\_min.z} defines the lowest permitted point.
Set the limits to match your actual machine and tooling before first use.
\end{warning}

\section{Seeding plates --- \code{[[plates]]}}
\label{sec:plates}

Each \code{[[plates]]} block describes one tray layout and appears as a selectable
button in the seeding workflow. Add as many blocks as you have tray types.

\begin{lstlisting}[language=TOML]
[[plates]]
name = "77pots"                          # button label in the UI
plate_size   = { x = 500, y = 700, z = 40 }  # tray dimensions in mm
first_pot    = { x = 7, y = 24 }         # centre of the first pot (robot coords, mm)
pot_distance = { x = 10, y = 12 }        # centre-to-centre spacing in mm
nb_pot       = { x = 7, y = 11 }         # grid size: pots in X and Y
\end{lstlisting}

The application computes the position of pot $(i,j)$, with
$0 \le i < \code{nb\_pot.x}$ and $0 \le j < \code{nb\_pot.y}$, as:
\[
  x_{ij} = \code{first\_pot.x} - i \cdot \code{pot\_distance.x},
  \qquad
  y_{ij} = \code{first\_pot.y} - j \cdot \code{pot\_distance.y}
\]

Note the \emph{subtraction}: the grid extends from the first pot towards
\textbf{negative} X and Y. The first pot is therefore the one with the largest
coordinates, and the whole grid must fit inside the safety limits of
Section~\ref{sec:limits}.

\begin{center}
\begin{tikzpicture}[scale=0.9]
  % Working area / axes
  \draw[-{Stealth}, thick, gray] (-6.8,0) -- (1.2,0) node[below right, black] {$X$};
  \draw[-{Stealth}, thick, gray] (0,-5.2) -- (0,1.0) node[above left, black] {$Y$};
  \node[below left, gray] at (0,0) {\small $0,0$};
  % Plate outline
  \draw[rounded corners=2pt, fill=accentlight, draw=accent, thick]
        (-6.3,-4.7) rectangle (0.8,0.6);
  \node[anchor=south west, accent] at (-6.25,-4.65) {\small tray};
  % Pot grid: first pot at (0.2, 0.1), spacing 1.2 / 1.15, 5 x 4 pots
  \foreach \i in {0,...,4}
    \foreach \j in {0,...,3}
      \draw[fill=white, draw=accent] ({0.2-\i*1.2},{0.1-\j*1.15}) circle (0.32);
  % First pot marker
  \draw[very thick, warncol] ({0.2-0.45},{0.1}) -- ({0.2+0.45},{0.1});
  \draw[very thick, warncol] ({0.2},{0.1-0.45}) -- ({0.2},{0.1+0.45});
  \node[warncol, anchor=west] at (1.0,0.9) {\small \code{first\_pot}};
  \draw[-{Stealth}, warncol] (1.35,0.75) to[bend right=15] (0.5,0.35);
  % Spacing annotations
  \draw[{Stealth}-{Stealth}, thick] ({0.2-1.2},{0.1+0.55}) -- ({0.2},{0.1+0.55})
        node[midway, above] {\small \code{pot\_distance.x}};
  \draw[{Stealth}-{Stealth}, thick] ({0.2+0.55},{0.1-1.15}) -- ({0.2+0.55},{0.1})
        node[midway, right] {\small \code{pot\_distance.y}};
  % Grid extent braces
  \draw[decorate, decoration={brace, mirror, amplitude=5pt}]
        ({0.2-4*1.2-0.35},-5.0) -- ({0.2+0.35},-5.0)
        node[midway, below=6pt] {\small \code{nb\_pot.x} columns};
  \draw[decorate, decoration={brace, amplitude=5pt}]
        (-6.6,{0.1-3*1.15-0.35}) -- (-6.6,{0.1+0.35})
        node[midway, left=6pt, rotate=90, anchor=south] {\small \code{nb\_pot.y} rows};
\end{tikzpicture}
\end{center}

\subsection{Worked example: adding a new tray type}

Suppose a new tray has $6 \times 9$ pots, 15\,mm apart in X, 14\,mm apart in Y, and
its first pot centre sits at robot coordinates $(30, 50)$:

\begin{lstlisting}[language=TOML]
[[plates]]
name = "54pots"
plate_size   = { x = 500, y = 700, z = 40 }
first_pot    = { x = 30, y = 50 }
pot_distance = { x = 15, y = 14 }
nb_pot       = { x = 6, y = 9 }
\end{lstlisting}

Check the extremes: the last pot lies at
$x = 30 - 5 \cdot 15 = -45$ and $y = 50 - 8 \cdot 14 = -62$ --- inside the default
$\pm 160$\,mm limits. Restart the application; the new \textbf{54pots} button appears
in the plate selection screen.
